\section{结论}
\label{sec:conclusion}

病例规模边际价值由估计目标与两类不可约方差共同决定。监测当期传播时规模收益不饱和，小规模下按检验功效确定资源。监测跨期均值时扩大单期规模面临收益递减；超过方差平衡规模后，新增资源应转向更多观测期、更广空间汇总、统一口径与更新基准模型。美国流感住院实证观察到与理论分解相容的收益递减模式，未被规模项解释的截距约占初期损失的一半，超出的部分主要不能由填报比例代理反映的报告波动解释；不过，实证尚不足以区分严格非零平台与缓慢衰减，也不能将截距唯一归因于环境方差。环境方差大小只能靠积累观测期认识，在季节性监测中这意味着积累流行季。允许偏差的平滑估计可进一步压低趋势误差，代价是对异常时期反应变慢，该权衡不在本文条件无偏框架内。
